\documentclass[a4paper,12pt]{article}
\usepackage[T2A]{fontenc}
\usepackage[utf8]{inputenc}
\usepackage[russian]{babel}
\usepackage{amsmath,amssymb,mathtools}
\renewcommand{\familydefault}{\sfdefault}
\usepackage{icomma}
\usepackage[a4paper,margin=20mm]{geometry}
\usepackage{microtype}
\usepackage{tikz}
\usetikzlibrary{calc,arrows.meta,positioning,shapes.geometric}
\usepackage{tabularx,booktabs,longtable}
\usepackage{enumitem}
\usepackage[hidelinks]{hyperref}
\usepackage{parskip}
\usepackage{fancyhdr}
\usepackage{setspace}
\usepackage{xcolor}

\definecolor{accent}{HTML}{2E6F70}
\definecolor{darkaccent}{HTML}{214F50}
\definecolor{textgray}{HTML}{2F3437}
\definecolor{softgray}{HTML}{6B7478}
\definecolor{linegray}{HTML}{D9DEDF}

\setlength{\headheight}{25pt}
\onehalfspacing
\setlength{\parindent}{0pt}
\setlength{\parskip}{0.55em}
\setlist[itemize]{leftmargin=1.45em,itemsep=0.25em,topsep=0.35em}
\setlist[enumerate]{leftmargin=1.75em,itemsep=0.55em,topsep=0.4em}
\renewcommand{\arraystretch}{1.32}

\pagestyle{fancy}
\fancyhf{}
\fancyhead[L]{\small\color{softgray}Ксения Клыкова}
\fancyhead[R]{\small\color{softgray}30.09.26}
\fancyfoot[C]{\small\color{softgray}\thepage}
\renewcommand{\headrulewidth}{0.3pt}
\renewcommand{\headrule}{\hbox to\headwidth{\color{linegray}\leaders\hrule height \headrulewidth\hfill}}

\newcommand{\sectiontitle}[1]{%
  \vspace{0.55em}%
  {\Large\bfseries\color{darkaccent}#1}\par
  \vspace{0.25em}\color{linegray}\hrule\color{textgray}\vspace{0.45em}
}
\newcommand{\subhead}[1]{{\bfseries\color{accent}#1}}
\newcommand{\checkline}[1]{\item[$\square$] #1}
\newcommand{\important}[1]{\textbf{\color{darkaccent}#1}}

\tikzset{
  flowstart/.style={ellipse, draw=accent, line width=0.9pt, text width=6.2cm,
    minimum height=1.0cm, align=center, fill=accent!5, font=\small},
  flowact/.style={rounded corners=4pt, rectangle, draw=accent, line width=0.9pt,
    text width=7.6cm, minimum height=1.1cm, align=center, fill=accent!3, font=\small},
  flowcond/.style={diamond, aspect=2.25, draw=darkaccent, line width=0.9pt,
    text width=4.8cm, align=center, fill=accent!4, inner sep=1.2pt, font=\small},
  flowarrow/.style={-{Stealth[length=3.2mm,width=2.1mm]}, line width=0.9pt, draw=darkaccent},
  flowlabel/.style={font=\small\bfseries, text=darkaccent, fill=white, inner sep=1.5pt}
}

\begin{document}
\color{textgray}

% ---------- Title page ----------
\thispagestyle{empty}
\begin{center}
\vspace*{2.4cm}
{\small\color{softgray}\MakeUppercase{математика}}\\[1.1cm]
{\Huge\bfseries\color{darkaccent}Чек-лист}\\[0.45cm]
{\LARGE\bfseries Поиск рациональных корней многочлена}\\[0.28cm]
{\Large Теорема Безу и схема Горнера}\\[1.35cm]

{\large Ксения Клыкова}\\[0.45cm]
{\normalsize 30 сентября 2026 г.}\\[1.5cm]

\begin{minipage}{0.72\textwidth}
\centering
\color{softgray}
Корень многочлена $\longleftrightarrow$ линейный множитель $\longleftrightarrow$ нулевой остаток.\\[0.35cm]
Схема Горнера позволяет быстро проверить кандидата и одновременно понизить степень многочлена.
\end{minipage}
\end{center}

\vfill
\begin{center}
{\small\color{darkaccent}Лёвин Артём Александрович эксклюзивно для Ксении Клыковой}
\end{center}
\vspace{0.6cm}
\begin{tikzpicture}[remember picture,overlay]
  \draw[accent!60,line width=1.15pt]
    ($(current page.south west)+(1.5cm,1.15cm)$)
    .. controls ($(current page.south west)+(6.0cm,2.0cm)$)
    and ($(current page.south east)+(-6.0cm,0.25cm)$)
    .. ($(current page.south east)+(-1.5cm,1.15cm)$);
\end{tikzpicture}
\clearpage

% ---------- Main checklist ----------
\sectiontitle{1. Теорема Безу: как узнать линейный множитель}

Для любого многочлена $P(x)$ при делении на $x-a$ остаток равен $P(a)$:
\[
P(x)=(x-a)Q(x)+P(a).
\]
Отсюда следует главный критерий:
\[
\boxed{\;P(a)=0\quad\Longleftrightarrow\quad x-a\text{ является делителем }P(x).\;}
\]

\begin{itemize}
  \checkline{Если Вы нашли число $a$ и получили $P(a)=0$, то $a$ --- корень уравнения $P(x)=0$.}
  \checkline{Одновременно можно вынести множитель $x-a$ и перейти к многочлену степени на единицу меньше.}
  \checkline{Если $P(a)\ne0$, число $a$ не является корнем; оно остаётся только проверенным кандидатом.}
\end{itemize}

\subhead{Пример.} Пусть
\[
P(x)=x^3-2x^2+3x-2.
\]
Проверяем $x=1$:
\[
P(1)=1-2+3-2=0.
\]
Значит, $x-1$ --- множитель. После деления получаем
\[
P(x)=(x-1)(x^2-x+2).
\]
Для квадратного множителя $D=1-8=-7<0$, поэтому над действительными числами исходное уравнение имеет единственный корень $x=1$.

\sectiontitle{2. Где искать рациональные корни}

Пусть
\[
P(x)=a_nx^n+a_{n-1}x^{n-1}+\dots+a_1x+a_0,
\]
где все коэффициенты целые. Если несократимая дробь
\[
r=\frac{p}{q},\qquad \gcd(p,q)=1,
\]
является рациональным корнем, то
\[
\boxed{p\mid a_0,\qquad q\mid a_n.}
\]

\begin{itemize}
  \checkline{Для приведённого многочлена ($a_n=1$) рациональный корень обязан быть целым делителем свободного члена $a_0$.}
  \checkline{Выписывайте делители со знаками $\pm$.}
  \checkline{Для неприведённого многочлена учитывайте знаменатели --- делители старшего коэффициента.}
  \checkline{Список кандидатов содержит возможные корни. Каждый кандидат всё равно нужно проверить.}
\end{itemize}

\subhead{Пример.} Для
\[
2x^3-3x^2-8x+12=0
\]
числитель рационального корня делит $12$, а знаменатель делит $2$. После сокращения кандидаты:
\[
\pm\frac12,\ \pm1,\ \pm\frac32,\ \pm2,\ \pm3,\ \pm4,\ \pm6,\ \pm12.
\]

\sectiontitle{3. Схема Горнера}

Схема Горнера одновременно выполняет две задачи:
\begin{itemize}
  \item вычисляет $P(a)$;
  \item если $P(a)=0$, даёт коэффициенты частного $P(x):(x-a)$.
\end{itemize}

\subhead{Правило.} Выпишите коэффициенты по убыванию степеней. Первый коэффициент снесите вниз. Затем повторяйте цикл
\[
\text{умножить на }a\;\longrightarrow\;\text{прибавить следующий коэффициент}.
\]
Последнее число --- остаток $P(a)$; остальные полученные числа --- коэффициенты частного.

\subhead{Пример.} Разделим
\[
P(x)=x^3-x^2+4x+6
\]
на $x+1=x-(-1)$.

\begin{center}
\begin{tabular}{c@{\hspace{1.1cm}}rrrr}
 & $1$ & $-1$ & $4$ & $6$\\
$-1$ &   & $-1$ & $2$ & $-6$\\
\midrule
 & $1$ & $-2$ & $6$ & $0$
\end{tabular}
\end{center}

Следовательно,
\[
P(x)=(x+1)(x^2-2x+6),
\]
а нулевой остаток подтверждает, что $x=-1$ --- корень.

\sectiontitle{4. Пропущенная степень означает нулевой коэффициент}

Перед схемой Горнера степени должны идти подряд. Если некоторой степени нет, впишите для неё коэффициент $0$.

\subhead{Пример.}
\[
x^3-3x+2=x^3+0x^2-3x+2,
\]
поэтому строка коэффициентов имеет вид
\[
1,\ 0,\ -3,\ 2.
\]
Для $a=1$ схема Горнера даёт
\[
1,\ 1,\ -2,\ 0,
\]
то есть
\[
x^3-3x+2=(x-1)(x^2+x-2)=(x-1)^2(x+2).
\]

\sectiontitle{5. Что важно контролировать}

\begin{itemize}
  \checkline{Не перепутать корень $a$ и делитель $x-a$: корню $-1$ соответствует множитель $x+1$.}
  \checkline{Не забыть нулевые коэффициенты при пропущенных степенях.}
  \checkline{Не считать кандидата корнем до проверки: нужен нулевой остаток.}
  \checkline{В схеме Горнера сначала умножаем, затем складываем.}
  \checkline{После понижения степени решить полученный многочлен до конца.}
  \checkline{Для дробного кандидата использовать несократимую дробь $p/q$ и условие $p\mid a_0$, $q\mid a_n$.}
  \checkline{Если после деления тот же корень снова обращает частное в ноль, корень кратный.}
\end{itemize}

\sectiontitle{6. Самопроверка}
\begin{itemize}
  \checkline{Я умею по значению $P(a)$ определить, является ли $a$ корнем.}
  \checkline{Я умею перейти от корня $a$ к множителю $x-a$.}
  \checkline{Я умею выписать целые и рациональные кандидаты на корни.}
  \checkline{Я умею выполнить схему Горнера без пропуска нулевых коэффициентов.}
  \checkline{Я умею восстановить частное по строке Горнера и решить оставшееся уравнение.}
\end{itemize}

% ---------- Flowchart ----------
\clearpage
\thispagestyle{plain}
\begin{center}
{\LARGE\bfseries\color{darkaccent}Алгоритм: поиск рационального корня и понижение степени}
\end{center}
\vspace{0.8cm}

\begin{center}
\begin{tikzpicture}[node distance=10mm and 24mm]
  \node[flowstart] (start) {Дано уравнение $P(x)=0$ с целыми коэффициентами};
  \node[flowact, below=of start] (cands) {Выпишите кандидатов $r=p/q$:\\ $p\mid a_0$, $q\mid a_n$, дробь $p/q$ несократима};
  \node[flowact, below=of cands] (pick) {Возьмите очередного кандидата $r$ и выполните схему Горнера};
  \node[flowcond, below=13mm of pick] (zero) {Последний элемент равен $0$?};
  \node[flowact, below=14mm of zero] (factor) {Да: $r$ --- корень, $x-r$ --- множитель.\\ Коэффициенты слева дают частное $Q(x)$};
  \node[flowact, below=of factor] (solve) {Решите $Q(x)=0$. Если степень всё ещё высокая, повторите алгоритм};
  \node[flowstart, below=of solve] (end) {Запишите полный набор корней};
  \node[flowact, right=30mm of zero, text width=5.0cm] (next) {Нет: этот кандидат не подходит. Возьмите следующий};

  \draw[flowarrow] (start) -- (cands);
  \draw[flowarrow] (cands) -- (pick);
  \draw[flowarrow] (pick) -- (zero);
  \draw[flowarrow] (zero) -- node[flowlabel,left]{да} (factor);
  \draw[flowarrow] (factor) -- (solve);
  \draw[flowarrow] (solve) -- (end);
  \draw[flowarrow] (zero.east) -- node[flowlabel,above]{нет} (next.west);
  \draw[flowarrow] (next.north) |- ($(pick.east)+(0.5cm,0)$) -- (pick.east);
\end{tikzpicture}
\end{center}
\clearpage

% ---------- Training block ----------
\sectiontitle{Тренировочный блок --- Путь А}
\begingroup
\small
\setstretch{1.14}
\textbf{10 задач.} Решения в этом разделе не приведены. Используйте теорему Безу, список рациональных кандидатов и схему Горнера.

\begin{enumerate}[label=\textbf{\arabic*.},itemsep=0.28em,topsep=0.25em]
  \item Для многочлена
  \[
  P(x)=x^3-5x^2+2x+8
  \]
  выпишите все возможные целые корни.

  \item Проверьте, является ли число $2$ корнем многочлена
  \[
  P(x)=x^3-3x^2-4x+12.
  \]

  \item Известно, что $P(-3)=0$. Какой линейный множитель обязательно входит в разложение $P(x)$?

  \item С помощью схемы Горнера разделите
  \[
  x^3-x^2+4x+6
  \]
  на $x+1$. Укажите частное и остаток.

  \item Для многочлена
  \[
  P(x)=x^4-5x^2+4
  \]
  запишите полную строку коэффициентов, которую нужно использовать в схеме Горнера.

  \item Решите уравнение
  \[
  x^3-4x^2+x+6=0,
  \]
  найдя сначала целый корень, а затем понизив степень схемой Горнера.

  \item Решите уравнение
  \[
  x^3-3x+2=0.
  \]
  Обратите внимание на пропущенную степень и кратность одного из корней.

  \item Для уравнения
  \[
  2x^3-3x^2-8x+12=0
  \]
  выпишите все различные рациональные кандидаты на корни в сокращённом виде.

  \item Решите уравнение
  \[
  2x^3+x^2-8x-4=0.
  \]
  Среди корней есть дробное число.

  \item Решите уравнение
  \[
  x^4-5x^2+4=0,
  \]
  последовательно понижая степень схемой Горнера после нахождения каждого подходящего целого корня.
\end{enumerate}
\endgroup

\clearpage

% ---------- Answers ----------
\sectiontitle{Ответы}
\begin{longtable}{@{}p{0.08\linewidth}p{0.84\linewidth}@{}}
\toprule
\textbf{№} & \textbf{Ответ}\\
\midrule
\endhead
1 & $\pm1,\ \pm2,\ \pm4,\ \pm8$.\\
2 & Да, так как $P(2)=0$.\\
3 & $x+3$.\\
4 & Частное $x^2-2x+6$, остаток $0$.\\
5 & $1,\ 0,\ -5,\ 0,\ 4$.\\
6 & $x=-1,\ 2,\ 3$.\\
7 & $x=-2$ и $x=1$; корень $x=1$ имеет кратность $2$.\\
8 & $\pm\frac12,\ \pm1,\ \pm\frac32,\ \pm2,\ \pm3,\ \pm4,\ \pm6,\ \pm12$.\\
9 & $x=-2,\ -\frac12,\ 2$.\\
10 & $x=-2,\ -1,\ 1,\ 2$.\\
\bottomrule
\end{longtable}

\vfill
\begin{center}
{\small\color{softgray}Лёвин Артём Александрович эксклюзивно для Ксении Клыковой}
\end{center}

\end{document}
